\subsection{Unconditional numerical selection on the original two-sign host}
\label{sec:unconditional-qs}

The averaged arithmetic comparison uses numerical companion and thin-sample
estimates, not a fixed-width zero-free strip for the auxiliary characters
at every height.  The following replacement retains both original
auxiliary primes, their size ranges, and their prescribed splitting
congruences.  It uses the classical prime number theorem in progressions
with the exceptional term retained, rather than uniform quasi-RH.
It was obtained by transferring the fixed-field Hecke companion
selection back to the rational case.

\begin{theorem}\label{thm:unconditional-qs}
Fix an integer \(M\ge1\), constants \(A,J\ge1\), and
\(0<\epsilon<1/52\).  For all sufficiently large \(U\), put
\[
 z=U^{\epsilon/5},\quad z'=z(\log U)^2,\quad Q=e^z,\quad D=e^{z'}.
\]
There are primes \(Q/2<q<Q\), \(D/2<d<D\), prime to \(2M\),
with
\[
 q\equiv3\bmod8,\qquad d\equiv7\bmod16,\qquad
 d\equiv-1\bmod q,\qquad d\equiv-1\bmod p
       \quad(3\le p\le z,\ p\text{ prime}).
\]
Let \(\xi\) run over the primitive characters with conductor dividing
\(8M\), and let \(\chi_l=(\,\cdot\,/l)\).
Use the primitive realization of every product of characters. For a
nonprincipal character define the corrected ramp
\[
 G_\chi(y)=\sum_{n\le y}\Lambda(n)\chi(n)(y/n-1)
                  +y\frac{L'}{L}(1,\chi),
\]
and use the principal correction \(y(\log y-1-\gamma)\) subtracted
from its ramp.  Then for \(U\le y\le U^J\), \(i,j\in\{0,1\}\),
\(i+j>0\),
\begin{equation}\label{eq:unconditional-companion}
 |G_{\xi\chi_q^i\chi_d^j}(y)|\le y^{1/2+\epsilon/2},\qquad
 \left|\frac{L'}{L}(1,\xi\chi_q^i\chi_d^j)\right|\le4\log y .
\end{equation}
The choices can be made so that all primitive nontrivial quadratic
companions of the same base characters multiplied by real dyadic
cyclotomic characters of height at most \(A(1+\log U)\) have a
zero-free box \(\Re s\ge1-\delta\), \(|\Im s|\le Q^2\) or \(D^2\)
according as \(j=0\) or \(j=1\), with fixed \(\delta>0\).
For the entire resulting family, including the untwisted characters,
\begin{equation}\label{eq:unconditional-thin}
 |G_\chi(v)|/v\le e^{-x^{\epsilon/8}}
 \quad (U\le x\le U^J,\ v\ge e^{x^{\epsilon/3}}/3).
\end{equation}
All thresholds depend only on the fixed data.  The assertion does
not include a fixed-width zero-free strip at unbounded height.
\end{theorem}

\begin{proof}
We first establish weighted correlations for selecting the two primes.
At the \(q\)-stage take \(V=Q\), a fixed modulus \(m_0\) containing
\(8M\), and a reduced class \(a\equiv3\bmod8\).
At the \(d\)-stage, after choosing \(q\), take \(V=D\) and the
modulus \(16q\prod_{3\le p\le z}p\) (prime factors only), enlarged
by fixed factors needed to avoid \(M\).  The prescribed class
\(a\equiv7\bmod16\), \(a\equiv-1\) at the other prescribed primes
is reduced. Denote the stage modulus by \(m\); at the second
stage \(\log m=O(z)\).  Fix any cutoff exponent \(L\) and moment
order \(2h\). Consider refined moduli \(ms\) where \(s\) is odd,
squarefree, prime to \(m\) and \(s\le(2U^L)^{2h}\).

Thorner--Zaman's relative prime number theorem for progressions
(\emph{Refinements to the prime number theorem for arithmetic
progressions}, Math.\ Z.\ 306 (2024), article 54,
\href{https://arxiv.org/abs/2108.10878}{Theorem 1.1})
applied to \((V/2,V]\) gives the logarithmically weighted count
in each reduced class \(b\) modulo \(ms\). Its main term is
\(\varphi(ms)^{-1}\int_{V/2}^V(1-\chi_*(b)t^{\beta_*-1})\,dt\)
if there is an exceptional real character, and the integrand is one
otherwise. Its relative error is bounded by
\[
 O\left(\exp\left[-c\frac{\log V}
 {1+\log(ms)+(\log V)^{2/5}(\log\log V)^{1/5}}\right]\right).
\]
The theorem's lower-main-term hypothesis holds even in the suppressed
class: the effective bound \(1-\beta_*\gg(ms)^{-1/2}\) and its
(1.7) give a main term at least \(V\exp(-O(1+\log(ms)))\), while
\(\log(ms)=o(\log V)\). Thus the main term eventually exceeds
\(V^{32/37+\eta_0}\) for a fixed small \(\eta_0>0\).
At the first stage all \(\log(ms)=O(\log U)\); at the second
they are \(O(z)\). The relative error is \(O(e^{-c'(\log U)^2})\)
at both stages.

By the Landau--Page theorem, up to the largest of these moduli there
is at most one primitive real character with a zero closer to one
than \(c_0/\log(3\max ms)\), with \(c_0\) sufficiently small.
Any other exceptional term is absorbed in the same
\(O(e^{-c'(\log U)^2})\) error relative to the unperturbed main term.
The single retained character, whenever induced modulo \(ms\), is
exceptional in the displayed theorem.  Thus the same exception is
used in all refinements. Give the candidate primes \(l\) weight
\(\log l\) in the fixed progression modulo \(m\).
Summing over the reduced classes modulo \(s\) shows that
\[
 \left|\mathbb E(l/s)\right|\ll e^{-c'(\log U)^2}
\]
for every \(s>1\) as above, apart from at most one \(s_*\).
Indeed if the exceptional conductor divides \(m\), its term is
constant across the refinement and cancels in the signed numerator.
The errors are relative to the same positive main terms and remain
small even when the original progression is suppressed. If the
conductor divides \(ms\) but not \(m\), its nontrivial character on
the \(s\)-factor pairs to zero unless it is precisely the primitive
quadratic character \((\,\cdot\,/s)\). This requires \(s\) to be the
part of its primitive conductor outside \(m\), so there is at most
one exception. Use the trivial correlation bound one for that
exception. The original progression has weight at least
\(V\exp(-O(1+\log m))\); consequently deleting \(O(V^\nu)\)
candidates, for any fixed \(\nu<1\), deletes a vanishing fraction
of its weight.

This gives a high-moment selection, rather than only a second moment.
Before selecting the primes, supply at most \(U^{C_0}\) sequences
\(b_i(p)\), with \(|b_i(p)|\ll\log p\) for \(p\le2U^L\),
where \(C_0,L\) are fixed. Delete \(2\) and the primes dividing
\(m\); their total absolute contribution is \(O(1+\log m)=O(z)\)
throughout the windows. Quadratic reciprocity expresses the remaining
sign \(\chi_l(p)\) as \((l/p)\) times a fixed sign. Expand the \(2h\)-th
moment of the real sum up to \(t\), and use real and imaginary parts
for complex sequences. The terms in which every prime appears an
even number of times cost at most
\((2h-1)!!(\sum |b_i(p)|^2)^h\). The terms with odd squarefree product
other than \(s_*\) cost \(O(e^{-c'(\log U)^2}(\sum|b_i(p)|)^{2h})\).
For product \(s_*\), remove one occurrence of each of its \(r\) prime
factors, leaving only even multiplicities. If the term exists,
\(r\le2h\) is even. Counting pairings bounds their absolute total by
\[
 C_h\prod_{p\mid s_*}|b_i(p)|
       \left(\sum_{p\le t}|b_i(p)|^2\right)^{h-r/2}
 \le C_h\left(\sum_{p\le t}|b_i(p)|^2\right)^h .
\]
Thus the moment is \(O_h((t\log(2t))^h)\), by Chebyshev.
Markov's inequality and the union bound over the sequences and
integer cutoffs \(U\le t\le2U^L\), with \(h\) sufficiently large,
give a vanishing bad fraction for the bound
\(O(t^{1/2+\epsilon/4})\).  The forced-prime contribution is smaller
than this bound. First choose \(q\) for the base sequences, then choose
\(d\) for the doubled sequences \(b_i(p)\), \(b_i(p)\chi_q(p)\).
This works after the deletion of any \(O(V^\nu)\) candidates at
each stage.

We use that deletion freedom to select genuine zero-free boxes.
Jutila's density theorem (\emph{On Linnik's constant},
Math.\ Scand.\ 41 (1977), 45--62,
\href{https://www.mscand.dk/article/download/11701/9717}{Theorem 1,
(1.8)}) bounds the total zero count over primitive characters of
conductor at most \(H\), \(\Re s\ge1-2\delta\), \(|\Im s|\le T\)
by \(O((H^2T)^{6\delta})\) for \(0<2\delta\le1/5\), taking its
fixed auxiliary exponent equal to one. At either selection stage
the conductors of all indicated dyadic companions are at most
\(V^{C_1}\) for a fixed \(C_1\); use \(T=V^2\).
Each primitive companion is ramified at the new candidate prime,
so distinct candidates have disjoint sets of witnesses to failure.
For sufficiently small fixed \(\delta\) the number of bad candidates
is \(O(V^{1/3})\), allowed by the preceding selection.
There is no growing factor for the number of dyadic characters.

Apply the moment argument to \(\xi(p)\log p\) for the fixed base
family, with \(L>J\) so large that
\(L\delta/2>J+\epsilon/5+3\).
The higher prime powers contribute \(O(\sqrt t)\), by Chebyshev,
so the actual primitive companion satisfies
\[
 P_\chi(t):=\sum_{n\le t}\Lambda(n)\chi(n)
      =O(t^{1/2+\epsilon/4})\quad(U\le t\le2U^L).
\]
The coefficients at new ramified primes are exactly zero, as
required. Let \(X=U^L\), and choose a smooth cutoff \(w\) equal
to one on \([0,1]\), zero on \([2,\infty)\).
The selected zero-free box and the standard Hadamard local
partial-fraction formula give
\(-L'/L(s,\chi)=O_\delta(\log(C_\chi(3+|t|)))\) in the smaller
strip \(\Re s\ge1-\delta/2\), \(|t|\le V^2/2\), where
\(C_\chi\) is the conductor. Shift the Mellin integral for
\(\sum \Lambda(n)\chi(n)w(n/X)/n\) from \(\Re w=1\) to
\(-\delta/2\), using the height \(V^2/2\).  The Mellin transform of
the smooth cutoff has residue one at zero and rapid vertical decay.
The vertical error is \(O(X^{-\delta/2}(1+\log C_\chi))\);
horizontal and tail errors are negligible by that rapid decay.
Since \(\log C_\chi=O(z')\), our choice of \(L\) gives
\[
 -L'/L(1,\chi)=\sum_n\frac{\Lambda(n)\chi(n)}n w(n/X)
                  +O(U^{-J-1}).
\]
This is the genuine logarithmic derivative, not a fitted
linear correction.  Substitution gives
\(G_\chi(y)=-P_\chi(y)-y\sum_{n>y}\Lambda(n)\chi(n)w(n/X)/n
+O(yU^{-J-1})\).
Partial summation yields \(O(y^{1/2+\epsilon/4})\) on the window.
Splitting the cutoff sum at \(U\) bounds the derivative by
\(\log U+O(1)+O(U^{-1/2+\epsilon/4})\) using the elementary
Mertens sum \(\sum_{n\le U}\Lambda(n)/n=\log U+O(1)\).
Absorbing constants proves \eqref{eq:unconditional-companion}.

Finally the standard classical Dirichlet zero-free region can be used
for the entire growing dyadic family. The selected boxes exclude
the possible real exceptional zero of each companion (shrink the
region constant to fit inside the selected width at height zero).
The untwisted real characters form a fixed finite family, since the
dual of the real dyadic tower has just two characters of order
dividing two. Their possible exceptional zeros are removed by
shrinking the fixed region constant once more. The principal pole
is allowed. Thus, with a fixed \(c>0\), the region has width
\(c/\log(C_\chi(3+|t|))\).
In the smoothed Perron formula for the corrected ramp at \(v\),
truncate the initial line \(1+1/\log v\) at
\(T=e^{\sqrt{\log v}}\) and move to
\(1-\min(1/8,c/(4\log(C_\chi(T+4))))\).
All nearby zeros are separated horizontally by the region, giving
a local partial-fraction bound \(O((1+\log(C_\chi(T+4)))^2)\).
The residue at one cancels the genuine correction, including the
double pole of the principal character. The initial tails are
\(O(v(1+\log v)/T)\), the horizontal edges have the quadratic
Perron denominator of size \(T^2\), and the remaining vertical
integral is bounded by its value of \(v^{\Re s}\) times a fixed
power of the logarithmic parameters. Consequently
\[
 \frac{|G_\chi(v)|}{v}\ll
 (1+\log C_\chi+\log v)^B
 \exp\left(-c'\frac{\log v}
                  {1+\log C_\chi+\sqrt{\log v}}\right)
\]
for fixed constants \(B,c'>0\).
Here the actual size bound is
\(\log C_\chi=O(x^{\epsilon/5}(\log(2x))^2)\),
uniformly for \(U\le x\le U^J\).
At \(\log v\ge x^{\epsilon/3}-\log3\), the exponent dominates
a fixed positive multiple of
\(\min(x^{2\epsilon/15}/(\log(2x))^2,x^{\epsilon/6})\),
which in turn dominates \(x^{\epsilon/8}\). This domination
also absorbs the prefactor uniformly along the entire tail:
until \(\log v\) reaches \(x^{2\epsilon/5}(\log(2x))^4\)
its logarithm is \(O(\log x)\), and beyond that the exponent
is bounded below by a constant times \(\sqrt{\log v}\).
This proves \eqref{eq:unconditional-thin}.
\end{proof}

\begin{remark}
The numerical estimates above remove the QRH premise of the
auxiliary analytic comparison on these actual scale windows.
They do not automatically supply a stronger all-height
zero-free condition used by a different source interface.
The geometric and full integral Selmer budgets remain separate.
The fixed-field Hecke version, using the same rational-prime
correlation argument and Thorner's Hecke zero-density theorem,
is proved in the repository's
\href{https://github.com/JoseBrox/GRH_distillation/blob/main/notes/unconditional-quadratic-selection.md}
{unconditional-selection note}.
\end{remark}
